\uuid{UmX2}
\titre{Application linéaire}
\auteur{hanani ; mzari}
\organisation{exo7}
\licence{CC BY-NC-SA 4.0}
\datecreate{2020-01-20}
\niveau{L1}
\module{Algèbre}
\chapitre{Application linéaire}
\sousChapitre{Définition}
\difficulte{3}
\theme{Applications linéaires, Application linéaire}
\exo7qcmid{lille-2.applications-lineaires.008}

\contenu{
\begin{question}
\qclassification{107.01}
\qtags{motcle=Application linéaire}
On considère les applications suivantes :
$$\begin{array}{rccc}f:&\Rr^3&\to&\Rr^2\\
& (x,y,z)&\to &(x-y,y+2z+a) \,  \end{array}  \quad \mbox{et} \quad  \begin{array}{rccc}g:&\Rr^3&\to&\Rr\\
& (x,y,z)&\to &(ax+b)(x+y).\end{array} $$
où $a$ et $b$  sont des réels. Quelles sont les assertions vraies ?
\begin{answers}
\bad{Pour tout $a\in \Rr$, $f$ est une application linéaire}.
\good{$f$ est une application linéaire si et seulement si $a=0$}.
\bad{$g$ est une application linéaire si et seulement si $a=b=0$}.
\good{$g$ est une application linéaire si et seulement si $a=0$}.
\end{answers}
\begin{explanations} Si $a\neq 0$, alors $f(0,0,0)=(0,a)\neq(0,0)$, donc $f$ n'est pas linéaire. On vérifie aussi que, si $a=0$, alors $f$ est linéaire.
\vskip0mm
Si $a\neq 0$, $g(2(1,0,0))=4a+2b$ et $2g(1,0,0)=2a+2b \neq 4a+2b$, donc $g$ n'est pas linéaire.
\vskip0mm
On vérifie que si $a=0$ et $b$ est quelconque, $g$ est linéaire.
\end{explanations}
\end{question}
}
